% !TeX root = ../../Book1.tex
% 纲要标题：### 16.2 可分次数
% 纲要位置：计划纲要.md，第 564 行。

\section{可分次数}
\label{b1:sec:1602}


% 纲要内容（仅作写作计划，尚非教材正文）：
%
% * 嵌入个数与可分次数
% * 可分次数的塔式公式
% * 纯不可分扩张及其正规性；在混合次数算例中计算最大可分中间域
%

计数嵌入时，生成元的像必须是对应最小多项式的根，并且逐层选择要保持此前已经确定的映射。上一章已给出延拓的存在性与次数上界；现在把重根造成的差额单独记下来。

\subsection{嵌入个数与可分次数}
\label{b1:subsec:1602:01}

\begin{definition}\label{b1:def:separable-degree}
设 \(L/K\) 为有限域扩张，\(\Omega\) 为包含 \(L\) 的 \(K\) 的代数闭包。把整数
\[
[L:K]_{\mathrm s}\coloneqq\#\operatorname{Hom}_K(L,\Omega)
\]
称为扩张 \(L/K\) 的\term[kefen-cishu]{可分次数}[separable degree]。代数闭包之间的 \(K\) 同构通过复合给出嵌入集合的双射，所以该数不依赖所选代数闭包。
\end{definition}
\symidx{\([L:K]_{\mathrm s}\)：可分次数}

对单扩张 \(L=K(a)\)，\cref{b1:lem:simple-embedding-extension}把每个 \(K\) 嵌入与 \(a\) 的一个共轭根对应起来，因而
\[
[K(a):K]_{\mathrm s}
=\#\{\,b\in\Omega\mid m_{a,K}(b)=0\,\}
\leq[K(a):K]=\deg m_{a,K}.
\]
总次数是向量空间维数，可分次数是不同嵌入的个数，而一个根出现多次并不会给出多个嵌入。在特征 \(p>0\) 时，若 \(m_{a,K}(x)=g(x^{p^r})\)，其中 \(g\) 不可约且可分，\cref{b1:prop:inseparable-polynomial-shape}给出
\[
[K(a):K]=p^r\deg g,\qquad
[K(a):K]_{\mathrm s}=\deg g,\qquad
\frac{[K(a):K]}{[K(a):K]_{\mathrm s}}=p^r.
\]
这里第三个数字正是各根共有的重数，记录总次数中没有转化为不同嵌入的因子。对一般有限扩张，这个比值也总是正整数；\cref{b1:prop:separable-inseparable-decomposition}将证明这一点，并把它定义为不可分次数。

\begin{theorem}\label{b1:thm:separable-embeddings}
有限扩张 \(L/K\) 满足 \([L:K]_{\mathrm s}\leq[L:K]\)，等号成立当且仅当 \(L/K\) 可分。若 \(L\) 由有限多个在 \(K\) 上可分的元素生成，则 \(L/K\) 可分。
\end{theorem}
\begin{proof}
固定包含 \(L\) 的 \(K\) 的代数闭包 \(\Omega\)。上界已由\cref{b1:prop:embedding-degree-bound}证明。假设 \(L=K(a_1,\ldots,a_r)\)，每个 \(a_i\) 在 \(K\) 上可分。令 \(E_0=K\)、\(E_i=K(a_1,\ldots,a_i)\)。每个 \(E_i\) 的嵌入都由 \(E_{i-1}\) 的一个嵌入延拓而来；若在同一层每个已有嵌入都有相同数目的延拓，则该层的嵌入总数就是原总数乘以这个数。下面要证明，第 \(i\) 层的延拓数恰为 \([E_i:E_{i-1}]\)。

记 \(a_i\) 在 \(E_{i-1}\) 上的最小多项式为 \(m_i\)。它整除 \(a_i\) 在 \(K\) 上的可分最小多项式，故仍可分。于是存在 \(u_i,v_i\in E_{i-1}[x]\)，使
\[
u_i m_i+v_i m_i'=1.
\]
给定 \(K\) 嵌入 \(\tau:E_{i-1}\rightarrow\Omega\)，逐项变换系数与形式求导相容，即 \(\tau(m_i')=\tau(m_i)'\)。把 \(\tau\) 施于上式的系数，得到
\[
\tau(u_i)\tau(m_i)+\tau(v_i)\tau(m_i)'=1.
\]
因此 \(\tau(m_i)\) 与其导数互素，仍没有重根；系数嵌入不把非零首项系数变成零，所以其次数仍为 \(\deg m_i=[E_i:E_{i-1}]\)。在代数闭包 \(\Omega\) 中，它恰有这么多个不同根。由\cref{b1:lem:simple-embedding-extension}，\(\tau\) 也恰有这么多个延拓。不同的限制对应互不相交的延拓集合，从 \(E_0=K\) 的唯一 \(K\) 嵌入开始逐层计数，得
\[
\#\operatorname{Hom}_K(L,\Omega)
=\prod_{i=1}^r[E_i:E_{i-1}]=[L:K].
\]
特别地，若 \(L/K\) 可分，取任何有限生成集即得到等号。

反过来，设嵌入总数已等于 \([L:K]\)。任取 \(a\in L\)，记其最小多项式的不同根数为 \(s\)。从 \(K(a)\) 到 \(\Omega\) 恰有 \(s\) 个 \(K\) 嵌入。对其中任意一个嵌入，把 \(L/K(a)\) 写成逐次单扩张，每层延拓的选择数至多为相应最小多项式的次数，所以其延拓总数至多为 \(\bigl[L:K(a)\bigr]\)。将所有 \(L\) 的嵌入按它们在 \(K(a)\) 上的限制分组，便得第一个不等号；第二个不等号则来自不同根数不超过多项式次数：
\[
[L:K]\leq s\bigl[L:K(a)\bigr]\leq\bigl[K(a):K\bigr]\,\bigl[L:K(a)\bigr]=[L:K].
\]
首项等于嵌入总数是此方向的假设，末项等于总次数是扩张次数的塔式公式。首尾相等迫使所有不等号均为等号，所以 \(s=\bigl[K(a):K\bigr]\)，即 \(a\) 的最小多项式没有重根。任意 \(a\) 都可分，故扩张可分；这也说明前面由可分生成元得到的等号足以推出整个扩张可分。
\end{proof}

可分扩张的定义要求每个元素都可分，这个定理却把有限生成情形的检查缩小到生成元：只需验证 \(a_1,\ldots,a_r\) 的最小多项式没有重根，就能保证 \(K(a_1,\ldots,a_r)\) 中的所有元素可分。特别是，从 \(a,b\) 各自的最小多项式看不出 \(a-b,ab\) 及非零 \(a\) 的逆元的可分性时，可以把它们放入共同的有限扩张 \(K(a,b)\)。这与\cref{b1:prop:algebraic-elements-subfield}用有限维数代替显式消元的做法相同。

\begin{proposition}\label{b1:prop:separable-subfield}
任意扩张 \(L/K\) 中在 \(K\) 上代数且可分的元素构成一个子域。
\end{proposition}
\begin{proof}
若 \(a,b\) 都可分，\cref{b1:thm:separable-embeddings}说明有限扩张 \(K(a,b)/K\) 可分。其元素 \(a-b,ab\) 以及 \(a\neq0\) 时的 \(a^{-1}\) 因而均可分。底域元素的最小多项式为一次，所以也包含在这个集合中。
\end{proof}

\subsection{可分次数的塔式公式}
\label{b1:subsec:1602:02}

\begin{theorem}\label{b1:thm:separable-degree-tower}
设 \(K\subseteq E\subseteq L\) 为代数扩张塔。则 \(L/K\) 可分当且仅当 \(E/K\) 与 \(L/E\) 都可分。若进一步有 \([L:K]<\infty\)，则可分次数满足
\[
[L:K]_{\mathrm s}=[L:E]_{\mathrm s}[E:K]_{\mathrm s}.
\]
\end{theorem}
\begin{proof}
先设 \([L:K]<\infty\)，固定包含 \(L\) 的 \(K\) 的代数闭包 \(\Omega\)；它也是 \(E\) 的代数闭包。把 \(L\) 的 \(K\) 嵌入按它们在 \(E\) 上的限制分组，一共有 \([E:K]_{\mathrm s}\) 种可能的限制。只要证明每组都恰有 \([L:E]_{\mathrm s}\) 个元素，便能相乘得到总数。因此考虑限制映射
\[
\operatorname{res}:\operatorname{Hom}_K(L,\Omega)
\longrightarrow\operatorname{Hom}_K(E,\Omega),
\qquad \sigma\longmapsto\sigma|_E.
\]
对每个 \(\tau\in\operatorname{Hom}_K(E,\Omega)\)，我们想把固定 \(E\) 的嵌入 \(\rho:L\rightarrow\Omega\) 变为在 \(E\) 上等于 \(\tau\) 的嵌入。\(\rho\) 的值在 \(\Omega\) 中，未必属于 \(E\)，所以只定义在 \(E\) 上的 \(\tau\) 不能直接与它复合。为此，把\cref{b1:cor:intermediate-embedding-extension}用于塔 \(K\subseteq E\subseteq\Omega\)，将 \(\tau\) 延拓为整个 \(\Omega\) 的一个 \(K\) 自同构 \(\widehat\tau\)。若 \(\rho\in\operatorname{Hom}_E(L,\Omega)\)，则 \(\widehat\tau\circ\rho\) 固定 \(K\)，并且对每个 \(e\in E\) 有
\[
(\widehat\tau\circ\rho)(e)=\widehat\tau(e)=\tau(e).
\]
所以它属于纤维 \(\operatorname{res}^{-1}(\tau)\)。反过来，若 \(\psi\in\operatorname{res}^{-1}(\tau)\)，则对每个 \(e\in E\) 有
\[
(\widehat\tau^{-1}\circ\psi)(e)
=\widehat\tau^{-1}(\tau(e))=e,
\]
故 \(\widehat\tau^{-1}\circ\psi\in\operatorname{Hom}_E(L,\Omega)\)。这两个复合互为逆映射，给出了 \(\operatorname{Hom}_E(L,\Omega)\) 与该纤维之间的双射。计数只要求存在一个双射，并不要求双射唯一：\(\widehat\tau\) 未必唯一，但任一选择都给出上述双射，足以说明纤维的大小等于同一个数 \([L:E]_{\mathrm s}\)。目标集合有 \([E:K]_{\mathrm s}\) 个元素，得到乘法公式。

有限情形下，每层可分次数至多等于该层次数。将刚证明的公式与\cref{b1:thm:tower-law}比较，两个有限正数的乘积取满值当且仅当每个因子取满值，再由\cref{b1:thm:separable-embeddings}得到可分性的塔判据。

一般情形，若 \(L/K\) 可分，限制到 \(E\) 以及扩大最小多项式的系数域，分别说明 \(E/K\)、\(L/E\) 可分。反过来，假设这两层都可分。即使 \(E/K\) 无限，一个具体多项式也只涉及有限多个系数，因而可以像\cref{b1:thm:algebraic-tower}的证明一样，把问题放入有限塔。给定 \(a\in L\)，写出其在 \(E\) 上的可分最小多项式
\[
f(x)=x^n+c_{n-1}x^{n-1}+\cdots+c_0,
\qquad E_0=K(c_0,\ldots,c_{n-1}).
\]
每个 \(c_i\) 都在 \(K\) 上代数且可分，所以\cref{b1:thm:separable-embeddings}说明 \(E_0/K\) 有限可分。

因为 \(f\) 在 \(E[x]\) 中不可约且可分，\(f'\neq0\)。其系数已经全在 \(E_0\) 中，缩小系数域不改变求导公式，故这个非零多项式 \(f'\) 也属于 \(E_0[x]\)。若 \(f\) 在 \(E_0[x]\) 中可约，同一分解会使它在 \(E[x]\) 中可约，矛盾；所以 \(f\) 在 \(E_0[x]\) 中仍不可约。再由\cref{b1:prop:irreducible-separable-derivative}，\(f\) 在 \(E_0[x]\) 中可分，故 \(E_0(a)/E_0\) 有限可分。应用已经证明的有限塔判据，\(E_0(a)/K\) 可分，于是 \(a\) 在 \(K\) 上可分。任意 \(a\) 均如此，结论成立。
\end{proof}
例如特征零的有限扩张中可分次数总等于次数；在 \(\mathbb F_p(t^{1/p})/\mathbb F_p(t)\) 中，可分次数为一而次数为 \(p\)。两者反映的是不同根的选择数与向量空间大小之间的差别。

\subsection{纯不可分扩张}
\label{b1:subsec:1602:03}

\begin{definition}\label{b1:def:purely-inseparable}
设 \(L/K\) 为代数扩张。如果 \(\operatorname{char}K=p>0\)，并且对每个 \(a\in L\) 都存在整数 \(r=r(a)\geq0\) 使 \(a^{p^r}\in K\)，则称 \(L/K\) 为\term[chun-bukefen-kuozhang]{纯不可分扩张}[purely inseparable extension]。在特征零情形，只把平凡扩张 \(K/K\) 称为纯不可分扩张。对任意在 \(K\) 上代数的元素 \(a\)，若 \(K(a)/K\) 纯不可分，就称 \(a\) 在 \(K\) 上纯不可分。
\end{definition}

这里的量词是对每个 \(a\in L\) 存在一个可以依赖 \(a\) 的指数 \(r=r(a)\)，没有要求一开始就用同一个指数处理所有元素。特征零时，不可约多项式全部可分；若它只有一个不同根，次数只能为一，根已经在 \(K\) 中。因此只有平凡扩张符合这种行为，这也解释了定义中的特征零约定。

\begin{proposition}\label{b1:prop:purely-inseparable-criteria}
设 \(K\) 为域，\(\Omega\) 为 \(K\) 的代数闭包，\(a\in\Omega\) 在 \(K\) 上代数。则 \(a\) 在 \(K\) 上纯不可分，当且仅当其最小多项式在 \(\Omega\) 中只有一个不同根。当这两个等价条件成立且 \(\operatorname{char}K=p>0\) 时，该最小多项式形如 \(x^{p^r}-c\)，其中 \(r\geq0\) 为整数、\(c\in K\)。因此，对任意代数扩张 \(L/K\)，固定一个包含 \(L\) 的 \(K\) 的代数闭包以后，扩张 \(L/K\) 纯不可分，当且仅当从 \(L\) 到该代数闭包的 \(K\) 嵌入只有一个。
\end{proposition}
\begin{proof}
先设 \(\operatorname{char}K=p>0\)。若 \(a^{p^r}=c\in K\)，则 \(x^{p^r}-c=(x-a)^{p^r}\)，其任何非常数因子都只有根 \(a\)。反过来，由\cref{b1:prop:inseparable-polynomial-shape}，只有一个不同根的不可约多项式的可分部分必须为一次，故它就是 \(x^{p^r}-c\)，并有 \(a^{p^r}=c\in K\)。这也保证整个 \(K(a)/K\) 纯不可分：其中每个元素都可写成 \(a\) 的 \(K\) 系数有理式，取 \(p^r\) 次幂后各系数与 \(a^{p^r}\) 都在 \(K\) 中，分母仍非零，因而所得值在 \(K\) 中。特征零时不可约多项式可分，只有一个根就只能为一次，等价于 \(a\in K\)。

固定自然包含 \(\iota:L\hookrightarrow\Omega\)。纯不可分时，每个 \(a\in L\) 的像必须是其最小多项式的根，而唯一的根就是 \(a\) 自身，故
\[
\operatorname{Hom}_K(L,\Omega)=\{\iota\}.
\]
若扩张不是纯不可分，某个 \(a\) 的最小多项式 \(m\) 有不同根 \(b\neq a\)。两次代入分别给出
\[
K(a)\cong_K K[x]/(m)\cong_K K(b),
\]
从而得到把 \(a\) 送到 \(b\) 的 \(K\) 嵌入 \(K(a)\rightarrow\Omega\)。再由\cref{b1:thm:embedding-extension}把这个嵌入延拓到 \(L\)，便得到不同于 \(\iota\) 的嵌入。
\end{proof}

特别地，若 \(\operatorname{char}K=p>0\)，且 \(m_{a,K}(x)=x^{p^r}-c\)，则在代数闭包中
\[
m_{a,K}(x)=(x-a)^{p^r},\qquad
[K(a):K]=p^r,\qquad [K(a):K]_{\mathrm s}=1.
\]
只有一个嵌入可选，总次数中的其余因子完全来自这一个根的重数。

\begin{corollary}\label{b1:cor:purely-inseparable-normal}
设 \(L/K\) 为纯不可分代数扩张，不要求有限。则 \(L/K\) 正规，且 \(\operatorname{Aut}_K(L)=\{\operatorname{id}_L\}\)。
\end{corollary}
\begin{proof}
固定包含 \(L\) 的 \(K\) 的代数闭包 \(\Omega\)。若 \(K[x]\) 中的不可约多项式在 \(L\) 中有根 \(a\)，其首一化就是 \(a\) 的最小多项式。由\cref{b1:prop:purely-inseparable-criteria}，它在 \(\Omega\) 中只有根 \(a\)，因此已经在 \(L[x]\) 中完全分裂。由正规扩张的定义，\(L/K\) 正规。每个 \(K\) 自同构同样必须把 \(a\in L\) 送到其最小多项式的根，所以固定每个元素，只能为 \(\operatorname{id}_L\)。
\end{proof}

\begin{proposition}\label{b1:prop:separable-inseparable-decomposition}
设 \(L/K\) 为有限域扩张，令 \(E\subseteq L\) 为其中所有在 \(K\) 上可分的元素组成的子域。则 \(E\) 是 \(L/K\) 的最大可分中间域：\(E/K\) 可分，且任何中间域 \(M\) 若 \(M/K\) 可分就有 \(M\subseteq E\)。此外，\(L/E\) 纯不可分，而且
\[
[L:K]_{\mathrm s}=[E:K].
\]
由扩张次数的塔式公式，商
\[
[L:K]_{\mathrm i}\coloneqq\frac{[L:K]}{[L:K]_{\mathrm s}}
=\frac{[L:K]}{[E:K]}=[L:E]
\]
是正整数，称为\term[bukefen-cishu]{不可分次数}[inseparable degree]；当 \(\operatorname{char}K=p>0\) 时它为 \(p\) 的幂，特征零时为一。
\end{proposition}
\begin{proof}
\cref{b1:prop:separable-subfield}保证 \(E\) 是域，它按定义对 \(K\) 可分。若中间域 \(M/K\) 可分，其每个元素都属于 \(E\)，所以 \(M\subseteq E\)。

正特征时，要让 \(a\in L\) 落到可分元素组成的域 \(E\) 中，可以从它的最小多项式出发。由\cref{b1:prop:inseparable-polynomial-shape}写成 \(m_{a,K}(x)=g(x^{p^r})\)，其中 \(g\) 不可约且可分。令 \(b=a^{p^r}\)，则 \(g(b)=0\)，所以 \(g\) 是 \(b\) 在 \(K\) 上的最小多项式，\(b\) 在 \(K\) 上可分。这正说明 \(a^{p^r}\in E\)，于是每个 \(a\) 在 \(E\) 上纯不可分，\(L/E\) 纯不可分。特征零时所有元素都可分，故 \(E=L\)。

由\cref{b1:prop:purely-inseparable-criteria}，\([L:E]_{\mathrm s}=1\)；又因 \(E/K\) 可分，\([E:K]_{\mathrm s}=[E:K]\)。可分次数塔式公式遂给出
\[
[L:K]_{\mathrm s}=[L:E]_{\mathrm s}[E:K]_{\mathrm s}=[E:K].
\]
再用总次数的塔式公式，商 \([L:K]/[L:K]_{\mathrm s}=[L:E]\) 为正整数。

若 \(\operatorname{char}K=p>0\)，取有限生成集 \(L=E(a_1,\ldots,a_m)\)，令 \(F_0=E\)、\(F_i=E(a_1,\ldots,a_i)\)。每个 \(a_i\) 的某个 \(p\) 次幂迭代已经在 \(E\subseteq F_{i-1}\) 中，所以 \(a_i\) 在 \(F_{i-1}\) 上纯不可分。由\cref{b1:prop:purely-inseparable-criteria}，其最小多项式形如 \(x^{p^{s_i}}-c_i\)，其中 \(s_i\geq0\)、\(c_i\in F_{i-1}\)。于是
\[
[F_i:F_{i-1}]=p^{s_i},\qquad
[L:E]=\prod_{i=1}^m p^{s_i}=p^{s_1+\cdots+s_m}.
\]
特征零时 \(E=L\)，不可分次数为一。
\end{proof}
\symidx{\([L:K]_{\mathrm i}\)：不可分次数}

上述结论把三种次数联系在同一域塔 \(K\subseteq E\subseteq L\) 中：下层 \(E/K\) 可分，其次数是 \([L:K]_{\mathrm s}\)；上层 \(L/E\) 纯不可分，其次数是 \([L:K]_{\mathrm i}\)，而
\[
[L:K]=[E:K][L:E]=[L:K]_{\mathrm s}[L:K]_{\mathrm i}.
\]
域 \(E\) 由 \(L\) 中所有 \(K\) 可分元素唯一确定。这里得到的是一层可分、上一层纯不可分的域塔，并没有一般声称还存在纯不可分中间域 \(P/K\)，使 \(L=EP\)。

\begin{corollary}\label{b1:cor:simple-extension-separable-part}
设 \(K\) 为特征 \(p>0\) 的域，\(L=K(a)\) 为其有限扩张。若 \(a\) 在 \(K\) 上的最小多项式为 \(m_{a,K}(x)=g(x^{p^r})\)，其中 \(r\geq0\) 为整数，\(g\in K[x]\) 首一、不可约且可分，则 \(L/K\) 的最大可分中间域为 \(E=K(a^{p^r})\)，并且
\[
[E:K]=\deg g,\qquad [L:E]=p^r.
\]
\end{corollary}
\begin{proof}
令 \(b=a^{p^r}\)。因为 \(g(b)=0\) 且 \(g\) 不可约，它就是 \(b\) 在 \(K\) 上的最小多项式。因此 \(K(b)/K\) 可分，且 \([K(b):K]=\deg g\)。由\cref{b1:prop:inseparable-polynomial-shape}，\(m_{a,K}\) 有 \(\deg g\) 个不同根；单扩张的嵌入对应于这些根，所以 \([L:K]_{\mathrm s}=\deg g\)。\cref{b1:prop:separable-inseparable-decomposition}中的最大可分中间域包含 \(K(b)\)，且次数同为 \(\deg g\)，故二者相等。最后由 \([L:K]=p^r\deg g\) 及扩张次数的塔式公式得到 \([L:E]=p^r\)。
\end{proof}

取奇素数 \(p\)，令 \(u\) 在 \(\mathbb F_p\) 上超越，并设 \(t=u^{2p}\)、\(K=\mathbb F_p(t)\)、\(L=\mathbb F_p(u)\)。多项式 \(x^{2p}-t\) 在 \(\mathbb F_p[t][x]\) 中对素元 \(t\) 满足 Eisenstein 判据，因而在 \(K[x]\) 中不可约，是 \(u\) 在 \(K\) 上的最小多项式。它又可写为
\[
x^{2p}-t=g(x^p),\qquad g(y)=y^2-t.
\]
\(g\) 同样满足 Eisenstein 判据，且因 \(p\) 为奇数而有 \(g'(y)=2y\neq0\)，所以不可约且可分。在代数闭包中，原多项式为 \((x-u)^p(x+u)^p\)，只有 \(u,-u\) 两个不同根，每个根重数为 \(p\)。

\ProseParagraphBreak

令 \(v=u^p\)，则 \(v^2=t\)，上述推论给出最大可分中间域
\[
E=K(v)=\mathbb F_p(u^p),\qquad
[E:K]=2,\qquad [L:E]=p.
\]
下层 \(E/K\) 的可分性来自最小多项式 \(y^2-t\) 的两个不同根；上层满足 \(u^p=v\)，是纯不可分扩张。于是 \([L:K]=2p\)、\([L:K]_{\mathrm s}=2\)、\([L:K]_{\mathrm i}=p\)：两个嵌入的选择来自下层，而上层的 \(p\) 倍次数全来自重数。
